\documentclass{beamer}
\usepackage{amsmath}
\usepackage{array}
\usepackage{tikz}
\usetikzlibrary{arrows.meta}
\usepackage{graphicx}
\usepackage{booktabs}
\usepackage{adjustbox}
\usepackage{multirow}
\usetheme{Madrid}
\newcounter{currentenumi}
\title [Energy Subsidy]{IMPACT OF FUEL SUBSIDY REMOVAL ON INFLATION AND CONSUMER EXPENDITURE IN NIGERIA}
%\subtitle{A Structural Vector Autoregression(SVAR) Approach}
\institute[Economist]{\textbf{Economist / Policy Researcher} \\
\vspace{8mm}
\url{iluahmad9@gmail.com} \\
\url{https://iluahmad.github.io} \\
\vspace{8mm}
\textbf{Working Paper Series: XXX03}}
\author{\textbf{Mr. Ahmad Ilu}}
\date{\today}
\begin{document}
  {
%\setbeamertemplate{background} 
%{
 
%}


\begin{frame}  
    \titlepage
\end{frame}

\begin{frame}{Outline}
    \tableofcontents
\end{frame}

\section{Introduction}

\begin{frame}{Introduction}
    \begin{itemize}
        \item Nigeria, the largest oil producer in Africa, discovered crude oil in 1956. which led to the establishment of its oil industry.
        \item Oil plays a critical role in Nigeria’s economy, serving as a major export and a key source of government revenue and foreign Receipts.
        \item Nigeria primarily produces light, sweet crude, valued for its low sulfur content. Approximately 37 billion barrels of proven oil reserves, ranking it among the top oil reserve holders globally. 
        \item Natural Gas Potential: Nigeria has substantial natural gas reserves (approximately 200 trillion cubic feet), positioning it as a key player in Africa’s gas market. Expanding natural gas production for domestic use and export (through pipelines and LNG) could diversify revenue sources and support economic growth.
        \item The Petroleum Industry Act (PIA): Enacted in 2021, the PIA aims to improve regulatory clarity, attract investment, and enhance the Nigerian National Petroleum Corporation’s (NNPC) operational efficiency. However, its impact is still unfolding as stakeholders adjust to the new regulations.
    \end{itemize}
\end{frame}
\section{Stylized Facts}

\begin{frame}{Stylized Facts}
 \resizebox{\textwidth}{!}
{
\begin{columns}[onlytextwidth]
    \column{0.7\textwidth}
    \begin{itemize}
        \item \textbf{Production}: Oil production fluctuates between 1.3 to 1.7 million barrels per day (bpd), impacted by market conditions, security concerns, and OPEC quotas.
        \item \textbf{Revenue Contribution}: Oil represents around 90\% of export revenues and 60-70\% of government revenue.
        \item \textbf{Industry Structure}: Dominated by international oil companies (Shell, Chevron, TotalEnergies) with increasing roles for NNPC.
        \item \textbf{Gas Reserves}: Holds around 200 trillion cubic feet of natural gas reserves.
        \item \textbf{Refineries}: Four government-owned refineries with a capacity of 445,000 barrels per day, and the privately-owned Dangote Refinery at 650,000 barrels per day.
    \end{itemize}
     \column{0.3\textwidth}
   \includegraphics[scale=0.285]{Rplot02.png} 
      \end{columns}
      }
\end{frame}

\section{Background of the Study}

\begin{frame}{Background of the Study}
  Subsidies, which represent the gap between consumer prices and efficient prices—or the practice of setting retail prices below fair market levels—are often intended to reduce inequality, alleviate energy poverty, and build political support. IMF (2023), asserted that Globally, fossil fuel subsidies were $\$7$ trillion in 2022, or 7.1 percent of GDP.  \\
 The reform of the fuel subsidy system is a critical component for redesigning the Nigerian economy and facilitating inclusive and sustainable economic diversification and growth. In recent years, fuel subsidies have consumed more than one-third of the recurrent budget, representing a substantial misallocation of resources that could be more effectively directed toward pro-poor initiatives. The apprehension regarding the political repercussions of significant price increases, alongside pervasive corruption and pressure from beneficiaries of the subsidy system, has led successive governments to hesitate in pursuing necessary reforms. 

\end{frame}

\begin{frame}{Subsidy Reform Attempts and Challenges}
    \begin{itemize}
        \item 2012: Nigeria attempted to remove the petrol subsidy, aiming to redirect the funds toward economic stability and infrastructure, but public protests forced a partial reinstatement.
        \item 2015: In 2015, as oil prices dropped, President Buhari’s administration introduced a phased subsidy removal plan, capping pump prices at NGN 145 per liter to stabilize costs. 
        \item 2020: Another attempt in 2020 to end fuel and electricity subsidies led to price hikes and public backlash, resulting in subsidy extensions through 2023.
       \item 2023: In 2023, President Bola Ahmed Tinubu’s administration finally removed the subsidy, calling it fiscally unsustainable. This led to significant fuel price increases, with costs rising incrementally from NGN 190 per liter to NGN 1050 per liter.
       \item Subsidies were draining government resources that could otherwise be allocated to critical sectors such as education, healthcare, infrastructure, and other social programs.
    \end{itemize}
\end{frame}
\begin{frame}{Economic Impacts of Fuel Subsidies}
       \textbf{2021}: PMS subsidies cost NGN 1.894 trillion, or 38\% of oil revenues.\textbf{2022}: Subsidy expenses increased to NGN 4.611 trillion, or 61.4\% of oil revenues.\textbf{2023}: Cost was NGN 3.135 trillion, or 33\% of oil revenues. \\
 However,the removal of PMS subsidy has significant political and economic implications for welfare and socioeconomic indicators. This policy change has contributed to persistent inflation, impacting the costs, pricing structures, and profit margins across various goods and services. The current trajectory of the general rise in prices is the pass-through effect from a compendium of politico-economic factors. The knock-on effect of Fuel, energy subsidy removal (electricity tariff hike), Currency redesign exercise, Sustained Naira depreciation, insecurity, flooding(Environmental), Climate Change, and adverse effects of beggar-thy-neighbor policy (border closure) have intensified poverty levels across board. This development has triggered an elusive shift in expenditure patterns amongst households, a shift characterized by a hand-to-mouth and out-of-pocket consumer economy.
\end{frame}
\section{Literature Review}
\begin{frame}{Literature Review: Theoretical Framework}
Theory of Subsidy and Welfare by Pigou (1920): Explores government interventions (subsidies and taxes) to address market inefficiencies and improve social welfare.
\begin{itemize}
\item Subsidies as Welfare Tools:
Subsidies can promote beneficial activities that would be under-produced or under-consumed in a free market.
Examples include subsidies for education and healthcare, which provide private benefits and positive spillovers to society (e.g., improved public health, productivity, and social cohesion).
Internalizing Positive Externalities: Subsidies help to internalize the positive externalities, compensating for the market's failure to reflect their full social value.
\item Pigovian Subsidy: A financial incentive designed to bridge the gap between private cost and social benefit, encouraging optimal consumption or production levels.
\item Pigou’s Tax Concept: Introduced taxes to address negative externalities, internalizing the costs of harmful effects and discouraging detrimental activities.
\end{itemize} 
\end{frame}
\begin{frame}{Literature Review: Compensating Variation (CV) and Equivalent Variation (EV)} 
    \begin{itemize} 
        \item \textbf{Compensating Variation (CV)}: Measures additional income needed for maintaining initial utility after a policy change. \\$CV = W(T_2, P_1) - W(T_1, P_1)$ \\
Where:
\begin{itemize}
    \item \( W(T_2, P_1) \) is the welfare (or utility) with the new policy or state (at time \(T_2\)) evaluated at the initial prices (\(P_1\)).
    \item \( W(T_1, P_1) \) is the welfare (or utility) in the baseline state (at time \(T_1\)) evaluated at the initial prices (\(P_1\)).
\end{itemize}
        \item \textbf{Equivalent Variation (EV)}: Quantifies money a consumer would pay to avoid a price change.\\
        $EV = W(T_1, P_1) - W(T_1, P_2)$ \\
Where:
\begin{itemize}
    \item \( W(T_1, P_1) \) is the welfare (or utility) in the baseline state (at time \(T_1\)) evaluated at the baseline prices (\(P_1\)).
    \item \( W(T_1, P_2) \) is the welfare (or utility) in the baseline state (at time \(T_1\)) evaluated at the new prices (\(P_2\)).
\end{itemize}
    \end{itemize}
\end{frame}
\begin{frame}{CV and EV: Measures of utility $\Delta$ Introduced by John Hicks (1939)}
    \begin{columns}[onlytextwidth]
    \column{0.5\textwidth}
    \begin{itemize}
    \item CV:  Refers to the amount of additional money an agent would need to reach their initial utility after a change in prices, or the introduction of new taxes or policy. 
    \item CV is forward-looking from the point before the change.
\item EV: It measures the amount of money a consumer would pay to avoid a price change, before it happens. 
\item EV is backward- looking from the point after the change.
 \end{itemize}
      \column{0.5\textwidth}
      \includegraphics[scale=0.36]{CV AND EV.png} 
      \end{columns}
\end{frame}
\begin{frame}{Theories of Consumption} 
The basic linear consumption function is defined as:
$C = C_0 + cY$
where:
\begin{itemize}
    \item \( C \) is the total consumption.
    \item \( C_0 \) is autonomous consumption (the level of consumption when income is zero).
    \item \( c \) is the marginal propensity to consume (MPC).
    \item \( Y \) is disposable income.
\end{itemize}
Theories of consumption are foundational in economics and aim to explain how individuals make choices about spending their income on goods and services over time. \\

\resizebox{\textwidth}{!}
{
%\begin{tabular}{lllp{8cm}}
%\begin{tabular}{| >{\centering\arraybackslash}m{3.5cm} | m{8cm} |}
\begin{tabular}{|c|c|} 
\hline
\textbf{Theory} & \textbf{Key Idea} \\
\hline
\textbf{Absolute Income Hypothesis} & Consumption rises with income, but less proportionately over time. \textit{(Keynes, 1936)} \\
\hline
\textbf{Relative Income Hypothesis} & Consumption influenced by income relative to others. \textit{(Duesenberry, 1949)} \\
\hline
\textbf{Life-Cycle Hypothesis} & Consumers plan to smooth spending over their lifetime. \textit{(Modigliani \& Brumberg, 1954)} \\
\hline
\textbf{Permanent Income Hypothesis} & Consumption based on expected average income. \textit{(Friedman, 1957)} \\
\hline
\textbf{Random Walk Hypothesis} & Consumption changes are unpredictable, following rational expectations. \textit{(Hall, 1978)} \\
\hline
\textbf{Behavioral Theories} & Psychological and social factors impact consumption choices. \textit{(Thaler \& Kahneman)} \\
\hline
\textbf{Precautionary Saving Model} & Saving as a safeguard against future uncertainty. \\
\hline
\textbf{Habit Formation Model} & Past consumption levels influence current spending. \\
\hline
\textbf{Intertemporal Choice Model} & Balances present vs. future consumption, factoring in time preferences. \textit{(Fisher)} \\
\hline
\end{tabular}
}
\end{frame}
\begin{frame}{Nigeria's Inflation Trajectory}
\begin{columns}[onlytextwidth]
    \column{0.7\textwidth}
 \textbf{Inflation Trends}:
Headline inflation rose to 33.20\% in August 2024, up from 25.79\% in August 2023; Core inflation reached 27.58\%, driven by higher input costs, insecurity, and infrastructure deficits; Food inflation climbed to 37.52\% from 29.34\% a year earlier, largely due to higher transportation costs from exchange rate unification and PMS subsidy removal.\\
\textbf{Signs of Moderation}: Both headline and food inflation have dropped for two consecutive months, attributed to the statistical base effect and seasonal variations.\\
\textbf{Future Inflation Risks}: Recent flooding in Northeastern states and PMS price adjustments by the NNPCL may counteract inflation deceleration, potentially limiting future price reductions.
    \column{0.3\textwidth}
      \includegraphics[scale=0.28]{PTRE.png}
      \end{columns}
\end{frame}
\begin{frame}{Pass-through Mechanism}
\begin{enumerate}
    \item \textbf{Exchange Rate Pass-Through to PMS Prices}
    \begin{itemize}
        \item \textbf{Exchange Rate Depreciation}: A depreciation of the naira increases the cost of importing goods, including Premium Motor Spirit (PMS), which Nigeria largely imports.
        \item \textbf{Higher PMS Import Costs}: The increased import cost for PMS raises domestic fuel prices to reflect the depreciated exchange rate.
        \item \textbf{Government Adjustments}: In the absence of subsidies, the government or NNPC adjusts PMS prices to cover the higher import costs, directly leading to increased fuel prices for consumers.
    \end{itemize}
    \item \textbf{PMS Price Increase to Overall Inflation}
    \begin{itemize}
        \item \textbf{Direct Effect on Transportation and Production Costs}: PMS is a key input for transportation and production. Higher PMS prices drive up transportation costs, impacting goods and services economy-wide.
        \item \textbf{Increased Cost of Goods}: Rising transportation costs increase prices for both raw and finished goods, contributing to \textbf{headline inflation}.
        \item \textbf{Impact on Food Inflation}: Food distribution logistics are sensitive to fuel prices. Higher PMS costs therefore significantly raise food prices, which are a major component of Nigeria's Consumer Price Index (CPI).
      
       \setcounter{currentenumi}{\theenumi}
    \end{itemize}
    \end{enumerate}
     \end{frame}
      
\begin{frame}
\begin{enumerate}
   \setcounter{enumi}{\thecurrentenumi}
   \item \textbf{Inflation to Reduced Real Income and Consumption}
    \begin{itemize}
        \item \textbf{Erosion of Real Income}: Rising inflation reduces consumers' \textbf{real disposable income}, meaning they can purchase fewer goods and services with the same nominal income.
        \item \textbf{Reduced Consumption}: Higher prices and diminished purchasing power lead households to cut back on discretionary and even essential spending, resulting in a decline in overall consumption.
        \item \textbf{Impact on Living Standards}: Lower consumption affects households' standard of living, potentially increasing poverty levels, especially among lower-income groups who are most vulnerable to price hikes.
    \end{itemize}

    \item \textbf{Feedback Loop and Long-Term Impact}
    \begin{itemize}
        \item \textbf{Persistent Inflation Expectations}: Sustained inflation may cause households to adjust their expectations, potentially leading to wage demands or more conservative spending patterns that can further entrench inflation.
        \item \textbf{Economic Slowdown}: Reduced consumption, alongside rising inflation, can stifle economic growth, creating a challenging environment for households and policymakers alike.
        
    \end{itemize}
\end{enumerate}    
       \end{frame}
       \begin{frame} {Pass-through FeedBack Loop}
       \includegraphics[scale=0.45]{Screenshot 2024-11-07 150749.png} 
       \end{frame}

\section{Methodology}
\begin{frame}{Wage indexation}
Wage indexation is the process of adjusting wages to account for changes in the cost of living, usually measured by an inflation index. The formula for wage indexation typically links wage increases to an inflation measure like the Consumer Price Index (CPI).\\ 
 \begin{exampleblock}{A common formula for Full wage indexation is:}
\begin{equation}
 W_t = W_{t-1} \times \left(1 + \frac{\Delta CPI}{CPI_{t-1}}\right)
\label{eqn:full indexation} 
\end{equation}
\end{exampleblock}
Where: \\
$ W_t$ = Wage at time t  \\
$ \Delta CPI$  = Change in the Consumer Price Index over the period \\
$ CPI_{t-1}$  = Consumer Price Index at the previous time period \\
There are several alternative formulas for wage indexing that can be used to adjust wages in response to inflation or other economic factors. \\
Here are a few examples outlined in the subsequent slides;
\end{frame}
\begin{frame}
\begin{block}{Partial Indexation Formula}
Partial indexation adjusts wages by a fraction of the CPI change, rather than the full amount. This can be useful in situations where employers want to share the burden of inflation with employees. \\
\begin{equation}
 W_t = W_{t-1} \times \left(1 + \alpha \times \frac{\Delta CPI}{CPI_{t-1}}\right)
\end{equation}
Where: \\
$\alpha$ = Fraction of the CPI change applied to wages (e.g. 0.5 / 50%)
    \end{block}
    \begin{exampleblock}{Capped Indexation Formula}
Capped indexation sets a maximum limit on the wage increase to prevent excessive adjustments during periods of high inflation. \\
\begin{equation}
W_t = W_{t-1}\times(1+ min \left(\frac{\Delta CPI}{CPI_{t-1}}, Cap\right)
\end{equation}
Where: \\
Cap = Maximum percentage increase allowed \\
    \end{exampleblock}
\end{frame}
\begin{frame}
\begin{block}{Sliding Scale Indexation Formula}
Sliding scale indexation adjusts wages based on different inflation rate brackets, providing more nuanced adjustments.
\begin{equation}
W_t = W_{t-1} \times \left(1 + \sum_{i=1}^{n} \alpha_i \times \frac{\Delta CPI_i}{CPI_{t-1}}\right)
\end{equation}
Where: $\alpha_{i}$ = Weight for each inflation bracket \textit{i} \\
 $\Delta CPI_{i}$ = CPI change within bracket \textit{i} \\
  \end{block}
  \begin{exampleblock}{Productivity-Linked Indexation Formula}
This formula incorporates productivity changes along with inflation to determine wage adjustments.
\begin{equation}
W_t = W_{t-1} \times \left(1 + \frac{\Delta CPI}{CPI_{t-1}} + \Delta P\right)
\end{equation}
Where: $\Delta_{P}$ = Change in productivity (e.g., percentage increase in output per worker)
    \end{exampleblock}
\end{frame}
\begin{frame} {Methodology}
\textbf{Data} \\
\begin{itemize}
\item Timeseries data for all five variables were obtained from the National Bureau of Statistics (NBS) and the Central Bank of Nigeria (CBN).
\item The study covers data from 2010 to 2024: Exchange Rate, Premium Motor Spirit (PMS) price, Inflation Rate, Consumer Price Index (CPI), and Final Household Consumption.
\item All series are originally in monthly series, except for Final Household Consumption data, which was interpolated to a monthly frequency from a quarterly series.
\item The selected period reflects the phase of full subsidy implementation without structural disruptions; 2016 onwards saw partial subsidy reductions and Price Modulation.
\item Previous studies( Kpodar \& Liu, 2021; Kpodar \& Abdallah, 2017) have highlighted inflation, PMS pump price, and exchange rate as key inflation drivers that significantly affect small businesses in emerging economies, such as Nigeria, following subsidy removal.
\end{itemize}
\end{frame}
\begin{frame}
    % Placeholder for methodology content
 A Unit Root Test is a statistical test used in time series analysis to determine whether a series is stationary or non-stationary. A stationary time series has statistical properties (like mean, variance) that do not change over time, while a non-stationary series does. Stationarity is crucial because non-stationary data can lead to spurious results in time series models, making it important to test for unit roots before proceeding with model estimations.
\begin{equation}
\gamma_t = \phi_1 + \rho \gamma_{t-1} + \mu_t \tag{1}
\end{equation}
Rearranging,
\begin{equation}
\gamma_t - \gamma_{t-1} = \phi_1 + \rho \gamma_{t-1} - \gamma_{t-1} + \mu_t \tag{2}
\end{equation}
\begin{equation}
\Delta \gamma_t = \phi_1 + (\rho - 1) \gamma_{t-1} + \mu_t \tag{3}
\end{equation}

Let \(\delta = \rho - 1\), then:

\begin{equation}
\Delta \gamma_t = \phi_1 + \delta \gamma_{t-1} + \mu_t \tag{4}
\end{equation}
\end{frame}
\begin{frame}
\textbf{Deterministic trend:}
   \[
   \Delta Y_t = \mu + \delta t + \rho Y_{t-1} + \sum_{i=1}^{p-1} \theta_i \Delta Y_{t-i} + \epsilon_t
   \]

 \textbf{Random Walk with drift:}
   \[
   \Delta Y_t = \mu + \rho Y_{t-1} + \sum_{i=1}^{p-1} \theta_i \Delta Y_{t-i} + \epsilon_t
   \]

 \textbf{ Random Walk without drift:}
   \[
   \Delta Y_t = \rho Y_{t-1} + \sum_{i=1}^{p-1} \theta_i \Delta Y_{t-i} + \epsilon_t
   \]
\end{frame}
\begin{frame}
\section*{The Brock, Dechert, and Scheinkman (BDS) Test}

The BDS test examines the null hypothesis that a time series is independently and identically distributed (i.i.d.). It is a non-parametric test based on spatial correlation, widely used to detect non-linear dependence in a series.

\subsection*{Steps of the BDS Test}
\textbf{BDS Statistic Calculation:}\\
   The BDS statistic is then calculated as:
   \[
   W_{m,\epsilon} = \frac{\sqrt{N} \left( C_m(\epsilon) - \left( C_1(\epsilon) \right)^m \right)}{\sigma_{m,\epsilon}}
   \]
   where \( \sigma_{m,\epsilon} \) is the standard deviation of \( C_m(\epsilon) - \left( C_1(\epsilon) \right)^m \).\\
\textbf{Decision Rule:} \\
   - If \( |W_{m,\epsilon}| \) is significantly large, we reject \( H_0 \), suggesting non-linear dependence.\\
   - Otherwise, we fail to reject \( H_0 \), implying the series could be i.i.d.\\
\subsection*{Interpretation}
- A significant BDS statistic indicates evidence of non-linear structures or dependencies in the series, which may suggest chaotic or complex patterns.
\end{frame}
\begin{frame}{Structural Vector Autoregression (SVAR)}
A Structural Vector Autoregression (SVAR) model is used to analyze the impact of different types of economic shocks by incorporating theoretical restrictions to identify causal relationships among variables.
\subsection*{VAR Representation}
A standard VAR model can be written as:
\begin{equation}
Y_t = A_1 Y_{t-1} + A_2 Y_{t-2} + \dots + A_p Y_{t-p} + u_t
\end{equation}
where:
\begin{itemize}
    \item \( Y_t \) is a vector of endogenous variables,
    \item \( A_i \) are coefficient matrices,
    \item \( u_t \) is a vector of residuals.
\end{itemize}
In the SVAR model, we introduce contemporaneous relationships among variables:
\begin{equation}
B Y_t = C_1 Y_{t-1} + C_2 Y_{t-2} + \dots + C_p Y_{t-p} + \epsilon_t
\end{equation}
where:
\begin{itemize}
    \item \( B \) is the matrix of structural coefficients,
    \item \( \epsilon_t \) is a vector of structural shocks (assumed to be uncorrelated).
\end{itemize}
\end{frame}
\begin{frame}
In Structural Vector Autoregressions (SVAR), Cholesky decomposition is used to identify structural shocks by decomposing the covariance matrix of the error terms. This method imposes a recursive ordering on the shocks.One of the approaches to recover the parameters in the structural form from the estimated parameters in the reduced form of equation is the use of recursive or Choleski factorization. This assumes Wold-chain ordering in which some variables cannot respond to other variables contemporaneously.
$$\varepsilon_t= B^{-1} \mu_t$$
\begin{equation}
\varepsilon_t = 
\begin{bmatrix}
\varepsilon_t^{Exchange Rate}\\
\varepsilon_t^{PMS}\\
\varepsilon_t^{Inflation}\\
\varepsilon_t^{Consumption Expenditure}
\end{bmatrix} =
\begin{bmatrix}
\alpha11 &0&0&0&0\\
\alpha21 &\alpha22&0&0&0\\
\alpha31 &\alpha32 &\alpha33&0&0\\
\alpha41 &\alpha42 &\alpha43&\alpha44&0\\
\end{bmatrix}
\begin{bmatrix}
\varepsilon_t^{Exch shock}\\
\varepsilon_t^{PMS shock}\\
\varepsilon_t^{\pi. shock}\\
\varepsilon_t^{Exp.Shock}\\
\end{bmatrix}
\end{equation}
When the elements of $B^{-1}$ are estimated, we can obtain the estimated vector of structural shocks, $\varepsilon_t$, since $\varepsilon_t= B^{-1}\mu_t$ alongside the responses of $x_t$ to each structural shocks in the system as demonstrated on equation.
\end{frame}
\begin{frame} {Method of Estimation}
\section*{Maximum Likelihood Estimation}
\textbf{{Maximum Likelihood Estimation}}\\
The Full Information Maximum Likelihood Estimation (MLE) is a method used to estimate the parameters of a statistical model by maximizing the likelihood function, which measures how likely it is to observe the given data under different parameter values. The likelihood function is derived from the joint distribution of the model's variables, taking into account the structural relationships between them.
The approximate log-likelihood for an SVAR(\(p\)) model is given by:
\begin{align*}
L(\theta) = & \ \text{cnst} + \frac{T - p + 1}{2} \left\{ \ln |A_0|^2 + \ln |\Omega_S^{-1}| \right\} \\
& - \frac{1}{2} \sum_{t = p + 1}^{T} \left( A_0 z_t - A_1 z_{t-1} - \dots - A_p z_{t-p} \right)' \Omega_S^{-1} \\
& \quad \left( A_0 z_t - A_1 z_{t-1} - \dots - A_p z_{t-p} \right).
\end{align*}

\end{frame}
\begin{frame} 
\section*{Bayesian Estimation for VAR Model}
To estimate the parameters of a VAR model using  \textbf{Bayesian methods}, we apply \textbf{Bayes' theorem}: which provides a framework for combining prior beliefs about the parameters with the likelihood of the observed data to obtain a posterior distribution.
\begin{equation}
p(\theta | Y) = \frac{p(Y | \theta) p(\theta)}{p(Y)}
\end{equation}

where:
\begin{itemize}
    \item \( p(\theta | Y) \) is the \textbf{posterior distribution} of the parameters given the data \( Y \),
    \item \( p(Y | \theta) \) is the \textbf{likelihood} of the data given parameters \( \theta \),
    \item \( p(\theta) \) is the \textbf{prior distribution} representing initial beliefs about \( \theta \),
    \item \( p(Y) \) is the \textbf{marginal likelihood} of the data, a normalizing constant.
\end{itemize}
To prepare forecasts, we need an estimate of \( \beta \). One approach is to use the mode of the posterior for \( \beta \), which can be found by maximizing:
\[
C(\beta) = \ln \{ L(Z | \beta, \Sigma_e) \} + \ln \{ p(\beta | \Sigma_e) \}
\]
where \( L(Z | \beta, \Sigma_e) \) represents the likelihood of the data \( Z \) given \( \beta \) and the error variance \( \Sigma_e \), and \( p(\beta | \Sigma_e) \) represents the prior distribution of \( \beta \) given \( \Sigma_e \).
\end{frame}
\section{Result Presentation and Discussion}
\begin{frame}{Descriptive Statistics}
    \centering
    \scriptsize % Adjust font size to fit table
    \begin{table}[h!]
        \centering
        \caption{Descriptive Statistics}
        \begin{adjustbox}{max width=\textwidth}
            \begin{tabular}{lcccc}
                \toprule
                \textbf{Statistic} & \textbf{Exchange Rate} & \textbf{PMS} & \textbf{Inflation Rate} & \textbf{Consumption} \\
                \midrule
                Mean & 332.3625 & 168.1171 & 14.4543 & 3567794 \\
                Median & 305.62 & 142 & 12.8 & 22119.92 \\
                Maximum & 1661 & 1030.54 & 34.19 & 41372991 \\
                Minimum & 150.08 & 65 & 7.7 & 7850.895 \\
                Std. Dev. & 272.417 & 155.4817 & 5.8682 & 10805160 \\
                Skewness & 3.0832 & 2.7780 & 1.4894 & 2.7967 \\
                Kurtosis & 13.4053 & 9.6825 & 5.2332 & 9.0001 \\
                Jarque-Bera & 1066.739 & 550.7029 & 101.0617 & 490.6394 \\
                Probability & 0 & 0 & 0 & 0 \\
                Sum & 58163.43 & 29420.49 & 2529.51 & $6.24 \times 10^8$ \\
                Sum Sq. Dev. & 12912715 & 4206371 & 5991.848 & $2.03 \times 10^{16}$ \\
                ADF & 3.1878** & -11.8999** & -5.2311** & -3.9600** \\
                PP & -12.6546** & -11.8999** & -5.2311** & -6.2400** \\
                Observations & 175 & 175 & 175 & 175 \\
                \bottomrule
            \end{tabular}
        \end{adjustbox}
    \end{table}
   Source: NBS,Author's Computation\\
    (**) Test Statistics at Significance at 5\%
\end{frame}

\begin{frame}{Result Presentation and Discussion}
\begin{columns}[onlytextwidth]
    \column{0.5\textwidth}
   \begin{table}[h!]
\centering
\caption{Wage Data with CPI and Real Wage Adjustments (2014–2024) using NBS Data}
\begin{adjustbox}{width=\textwidth}
\resizebox{\textwidth}{!}
{
\begin{tabular}{cccccc}
\toprule
\textbf{Year} & \textbf{CPI} & \textbf{Minimum Wage} & \textbf{Real Wage} & \textbf{Inflation Adjusted Minimum Wage} \\
\midrule
2014 & 157.4 & 18000 & 11435.4 & 18000.0 \\
2015 & 171.6 & 18000 & 1049.1 & 19621.2 \\
2016 & 198.3 & 18000 & 9077.0 & 22677.8 \\
2017 & 230.5 & 18000 & 7808.0 & 26363.2 \\
2018 & 257.3 & 18000 & 6996.0 & 29423.4 \\
2019 & 286.6 & 30000 & 10467.1 & 32776.6 \\
2020 & 322.2 & 30000 & 9312.0 & 36842.3 \\
2021 & 379.9 & 30000 & 7896.0 & 43449.3 \\
2022 & 447.2 & 30000 & 6707.9 & 51144.9 \\
2023 & 547.5 & 30000 & 5479.7 & 62608.2 \\
2024 & 733.4 & 30000 & 4090.8 & 83865.6 \\
\bottomrule
\end{tabular}
}
\end{adjustbox}
\end{table}
Using May 2014, as a base year, Nigeria inflation adjusted minimum wage in 2024 amounts to approximately ₦84,000 which is equivalent to about 56\$, at a rate of ₦1,500 / 1\$

  \column{0.5\textwidth}
    \begin{figure}[hbtp]
    \includegraphics[scale=0.35]{inflation adjusted minimum wage.png}
    \end{figure}
    
     \end{columns}
\end{frame}
\begin{frame}{IRFs}
\resizebox{\textwidth}{!}
{
\includegraphics[scale=1]{irf01.png} 
}
\end{frame}
\begin{frame}
\begin{columns}[onlytextwidth]

    \column{0.5\textwidth}
    \resizebox{\textwidth}{!}
{
    
    \includegraphics[scale=0.4]{fevd01.png} 
    }
    \column{0.5\textwidth}
    \resizebox{\textwidth}{!}
{
    \includegraphics[scale=0.4]{hedc01.png} 
    }
    \end{columns}
\end{frame}
\begin{frame} {Results Discussion}
The result presented from previous slides have shown the Descriptive Statistics from all model variables. it can be observed that all series are possibly skewed,and have leptokurtic distributions as all kurtosis test statistics are $>3$. In the same vein, it can be observed that both ADF and PP Test Statistics have indicated that all series have unit root at level, and stationary after taking the first difference, ie all are Integrated at order 1 I(1) at 5\% level of significance. \\
IRFs in line with economy theory were recursively used to analyze the dynamic impact of a one-time shock to one variable on itself and other variables in the system. In line with the postulations of this study, a one standard deviation positive exchange rate shock lead to an increase in exchange rate (₦/\$) to instantly depreciate further at period 2, persist up to period 10, This shock leads to an immediate increase of Price of PMS in the economy which pick up directly at period 2 through increases the cost of importing refined fuel and indirectly through landing costs that is reflected in domestic fuel prices.
\end{frame}
\begin{frame}
Subsequently, This steady rise in domestic prices feeds into general prices as cost-push shock, as PMS is a key input for transportation and production. Higher PMS prices drive up transportation costs, leading to hike in feedstock cost thereby impacting goods and services economy-wide and persistently driving headline, food inflation rates. \\
As observed from the IRFs above, we can see how abruptly and persistently the shock evolve and spikes into higher levels. Again, Food distribution logistics are sensitive to fuel prices. Higher PMS costs therefore significantly raise food prices, which are a major component of Nigeria's Consumer Price Index (CPI). 
Ultimately, this chain of shock transmission leads to a deliberate increase in household consumption expenditure, which disrupts consumption patterns, reduces savings, intensify vulnerabilities and heightens the demand for credit and social safety nets. \\
The Forecast Error Variance Decomposition (FEVD) analysis reveals that the exchange rate, inflation rate, and consumption expenditure exhibit strong exogeneity in both the short and long run, indicating that these variables are largely unaffected by other shocks in the system over time.
\end{frame}
 This suggests that shifts in exchange rates, inflation, and consumption expenditures are driven primarily by their own dynamics, rather than by external factors.\\
In contrast, the price of Premium Motor Spirit (PMS) or fuel is shown to be significantly influenced by fluctuations in consumption, particularly in the medium and long term. This dependency suggests that demand for fuel remains relatively inelastic; even as fuel prices change, consumption patterns remain steady, likely due to the essential nature of PMS for daily activities and economic productivity. Consequently, shifts in consumption expenditure directly impact fuel prices over time, reinforcing the notion that fuel demand is less responsive to price changes, which may reflect the lack of viable alternatives or substitutes in the market. \\
In essence, while core economic variables like exchange rate, inflation, and consumption hold steady against external disturbances, fuel prices are more vulnerable to consumption-driven shifts, underscoring a critical dependency within the Nigerian economy. This pattern points to potential challenges for policy interventions aimed at stabilizing fuel prices, as demand-side influences may dampen the effectiveness of such measures in both the medium and long term.

\begin{frame}{Scenarios Analysis}
This study established three (3) scenarios in the out of sample space and conditional forecast to empirically establish the path of fuel price in a simulation procedure to see the trajectory of inflation in Nigeria
for six (6)and twelve(12) months into the future given exchange rate movement.
\begin{itemize}
\item Baseline Scenarios:
This scenario assumes the continuation of current economic conditions without any major shocks to the exchange rate. It serves as a reference point to assess how fuel prices and inflation evolve under stable exchange rate conditions.
\item  50\% Appreciation of the Exchange Rate:
In this scenario, the exchange rate strengthens by 50\%, meaning the local currency gains value relative to foreign currencies. This appreciation is expected to reduce imported inflation and potentially lower fuel prices, which may ease inflationary pressures over the six-month forecast period.
\item 50\% Depreciation of the Exchange Rate:
This scenario simulates a weakening of the exchange rate, Such depreciation would likely increase the cost of imports, including fuel, thereby accelerating inflation. 
\end{itemize}
\end{frame}
\begin{frame}
 \begin{table}[ht]
\centering
\caption{Scenario Analysis Based on Exchange Rate Movement in Nigeria}
 \resizebox{\textwidth}{!}
{
\begin{tabular}{|c|c|c|c|c|}
\hline
\textbf{Scenario} & \textbf{Exchange Rate} & \textbf{PMS} & \textbf{Inflation Rate} & \textbf{Consumption} \\
\hline
\textbf{Baseline} & & & & \\
2025M03 & 1795.41 & 928.9535 & 36.03492 & 41999980 \\
2025M09 & 1973.138 & 1050.939 & 40.11438 & 50221910 \\
\hline
\textbf{Appreciation} & & & & \\
2025M03 & 897.705 & 794.2101 & 36.07536 & 55135410 \\
2025M09 & 986.569 & 968.6568 & 42.17932 & 70284870 \\
\hline
\textbf{Depreciation} & & & & \\
2025M03 & 2693.115 & 1063.697 & 35.99447 & 28864540 \\
2025M09 & 2959.707 & 1133.221 & 38.04944 & 30158960 \\
\hline
\end{tabular}
}

\end{table}
Analysis based on the table above, reveals that based on the 3 scenarios, That when Exchange rate Appreciate by 50\% in the next 12 months; Exchange Rate, PMS, Inflation Rate would amount to 986.569, 968.6568, and 42.17932 respectively. Once again,this implies the persistence of Nigerian inflation rate.
\end{frame}
\section{Policy Recommendations}

\begin{frame}{Policy Recommendations}
    % Placeholder for policy recommendations content
    Prudent utilization of funds from subsidy removal should focus on ensuring long-term economic stability and social welfare. This can be achieved through:
    \begin{itemize}
    \item Tailored safety nets should be put in place to protect the poor prior to the implementation of possible subsidy reforms
    \item A range of compensation mechanism could be implemented, tailored to the specific needs and capacity of each state and the Federal Capital Territory. These measures may include: Transport subsidies or vouchers, Public transportation initiatives, E-Wallets for smallholder farmers, Free school meal programs for children, Telemedicine for vulnerable groups, Alternative energy sources and stimulating of off-grid energy sources, and Vocational training and skills development programs
    \end{itemize}
\end{frame}
\begin{frame}
\begin{itemize}
\item Investing in critical infrastructure: Channelling the funds into the development of transportation, energy, and healthcare systems to improve productivity and quality of life.


\item Boosting economic diversification: Investing in sectors beyond oil, like agriculture, technology, and manufacturing, to reduce dependency on oil revenues and create sustainable growth.

\item Supporting job creation and skills development: Funding vocational training, entrepreneurship programs, and workforce development to improve employability and foster economic resilience.

\item Debt reduction and fiscal stability: Using a portion of the funds to reduce national debt or shore up government reserves, which can strengthen fiscal policy and enhance economic confidence.\\
These strategies would help ensure that the proceeds from subsidy removal contribute to long-term prosperity while minimizing short-medium term hardship.
\end{itemize}
\end{frame}
\begin{frame}
\begin{center}
\textbf{Scan to save my Contact} \\
\includegraphics[scale=0.39]{Barcode.jpg} 

\textbf{Thank you for Listening!} 
\end{center}
       
\end{frame}
\end{document}
